%% ----------------------------------------------------------------------------
%%  §1.1  Gli assiomi dei numeri reali
%% ----------------------------------------------------------------------------
\section{Gli assiomi dei numeri reali}
\label{sec:assiomi}

Il nostro punto di partenza è quello di \emph{assumere, come postulato, che esista il
sistema dei numeri reali}. Assumiamo cioè che esista un insieme, che chiamiamo insieme
dei numeri reali e che indichiamo con $\R$, i cui elementi si possono sommare e
moltiplicare e si possono confrontare tra loro, e che tali operazioni e tale confronto
obbediscano a un certo numero di regole. Sono le regole che il lettore ha sempre usato,
in modo naturale, da quando ha imparato a far di conto: per tutti è ovvio che $2+3$ sia
uguale a $3+2$, o che $2\cdot3=3\cdot2$. Esprimiamo questo fatto in generale, per due
numeri reali $a$ e $b$, dicendo che valgono le \emph{proprietà commutative} della
somma e del prodotto:
\begin{equation}\label{eq:commutativa}
  a+b=b+a, \qquad a\cdot b=b\cdot a .
\end{equation}
Il lettore conosce sicuramente (provi a scriverle da solo, prima di proseguire) anche
la proprietà associativa, sia rispetto alla somma sia rispetto al prodotto, e la
proprietà distributiva. Elenchiamo ora in modo completo le proprietà che assumiamo
valere \emph{per assioma}, dividendole in tre gruppi: le proprietà relative alle
operazioni, le proprietà relative all'ordinamento e l'assioma di completezza.

\begin{notastorica}[Che cosa sono i numeri reali?]
Per secoli i matematici hanno usato i numeri reali senza una definizione precisa: la
«retta dei numeri» era un'immagine geometrica sufficiente per Newton e per Eulero. Fu
solo nella seconda metà dell'Ottocento, con Richard Dedekind (1831--1916), Georg Cantor
(1845--1918) e Karl Weierstrass (1815--1897), che si sentì il bisogno di \emph{costruire}
$\R$ a partire dai numeri razionali, per dare un fondamento rigoroso al calcolo
infinitesimale. Dedekind, in \emph{Stetigkeit und irrationale Zahlen} (1872), racconta di
aver avuto l'idea delle sue «sezioni» il 24 novembre 1858, mentre preparava le lezioni
di calcolo al Politecnico di Zurigo, insoddisfatto di dover ricorrere all'intuizione
geometrica per dimostrare che una successione crescente e limitata ha un limite.
Nello stesso 1872 Cantor, e già nel 1869 il francese Charles Méray (1835--1911),
costruirono i reali per un'altra via, a partire dalle successioni di razionali che «si
stringono» (le successioni di Cauchy del \cref{sec:bw-cauchy}); Weierstrass ne esponeva
una sua nelle lezioni berlinesi. Il programma comune --- fondare l'intera Analisi sui
numeri naturali, attraverso i razionali --- fu chiamato da Felix Klein, nel 1895,
\emph{aritmetizzazione dell'analisi}. La frase di Kronecker in epigrafe ne riassume lo
spirito, anche se Kronecker stesso rifiutava le costruzioni di Dedekind e di Cantor,
che fanno uso di insiemi infiniti.
L'approccio di questo libro è diverso e più economico: non costruiamo i numeri reali,
ma \emph{postuliamo} che esistano con le proprietà elencate qui sotto. È l'approccio
\emph{assiomatico}, reso rigoroso da David Hilbert (1862--1943) nel 1900. I due approcci
sono equivalenti: si può dimostrare che esiste, a meno di «cambiare nome» agli elementi,
un solo sistema che soddisfa tutti gli assiomi che stiamo per enunciare (un modello
costruito con le sezioni di Dedekind è descritto nel \S\ref{subsec:modello-reali}).
\end{notastorica}

\subsection{Assiomi relativi alle operazioni}

Sono definite le operazioni di \emph{addizione} ($+$) e di \emph{moltiplicazione}
($\cdot$) tra coppie di numeri reali, con le seguenti proprietà ($a$, $b$, $c$ indicano
numeri reali generici).

\begin{assioma}[Proprietà delle operazioni]\label{ass:operazioni}
\leavevmode
\begin{enumerate}[label=\textup{(A\arabic*)}, ref=(A\arabic*), leftmargin=3em]
  \item\label{ass:associativa} \textbf{Proprietà associativa:}
    $(a+b)+c=a+(b+c)$ \ e \ $(a\cdot b)\cdot c=a\cdot(b\cdot c)$.
  \item\label{ass:commutativa} \textbf{Proprietà commutativa:}
    $a+b=b+a$ \ e \ $a\cdot b=b\cdot a$.
  \item\label{ass:distributiva} \textbf{Proprietà distributiva:}
    $a\cdot(b+c)=a\cdot b+a\cdot c$.
  \item\label{ass:neutri} \textbf{Esistenza degli elementi neutri:} esistono in $\R$ due
    numeri \emph{distinti}, indicati con $0$ e $1$, tali che $a+0=a$ e $a\cdot1=a$ per
    ogni $a\in\R$.
  \item\label{ass:opposto} \textbf{Esistenza degli opposti:} per ogni numero reale $a$
    esiste un numero reale, indicato con $-a$, tale che $a+(-a)=0$.
  \item\label{ass:inverso} \textbf{Esistenza degli inversi:} per ogni numero reale
    $a\neq0$ esiste un numero reale, indicato con $a^{-1}$, tale che $a\cdot a^{-1}=1$.
\end{enumerate}
\end{assioma}

\begin{osservazione}
Le proprietà (A1)--(A6) dicono, nel linguaggio dell'algebra, che $\R$ è un
\emph{campo}\index{campo}. Si noti che la sottrazione e la divisione non compaiono tra le
operazioni: esse sono \emph{definite} a partire dall'addizione e dalla moltiplicazione
ponendo
\[
  a-b \coloneqq a+(-b), \qquad \frac{a}{b}\coloneqq a\cdot b^{-1}\quad(b\neq0).
\]
Il simbolo $\coloneqq$ si legge «è definito uguale a». Si noti inoltre che nell'assioma
\ref{ass:inverso} si richiede $a\neq0$: vedremo nel \cref{sec:conseguenze} che questa
restrizione non è un capriccio, ma una necessità logica.
\end{osservazione}

\subsection{Assiomi relativi all'ordinamento}

È definita la relazione di \emph{minore o uguale} ($\le$) tra coppie di numeri reali,
con le seguenti proprietà.

\begin{assioma}[Proprietà dell'ordinamento]\label{ass:ordine}
\leavevmode
\begin{enumerate}[label=\textup{(O\arabic*)}, ref=(O\arabic*), leftmargin=3em]
  \item\label{ass:dicotomia} \textbf{Dicotomia:} per ogni coppia di numeri reali $a$, $b$ si
    ha $a\le b$ oppure $b\le a$.
  \item\label{ass:antisimmetrica} \textbf{Proprietà antisimmetrica:} se valgono
    contemporaneamente $a\le b$ e $b\le a$, allora $a=b$.
  \item\label{ass:ordine-somma} \textbf{Compatibilità con la somma:} se $a\le b$ allora
    $a+c\le b+c$ per ogni $c\in\R$.
  \item\label{ass:ordine-prodotto} \textbf{Compatibilità con il prodotto:} se $0\le a$ e
    $0\le b$ allora $0\le a+b$ e $0\le a\cdot b$.
\end{enumerate}
\end{assioma}

\begin{osservazione}
La proprietà di dicotomia \ref{ass:dicotomia}, applicata con $b=a$, dice in particolare
che $a\le a$ per ogni $a$ (\emph{proprietà riflessiva}). Tra gli assiomi non abbiamo
invece elencato la proprietà \emph{transitiva} («se $a\le b$ e $b\le c$ allora
$a\le c$»): la dimostreremo nel prossimo paragrafo come conseguenza degli altri assiomi.
Le relazioni $\ge$, $<$ e $>$ sono definite a partire da $\le$:
\begin{gather*}
  a\ge b \sse b\le a, \\
  a<b \sse (a\le b \text{ e } a\neq b), \qquad
  a>b \sse (a\ge b \text{ e } a\neq b).
\end{gather*}
Le relazioni $<$ e $>$ sono dette di \emph{minore stretto} e \emph{maggiore stretto}.
\end{osservazione}

\subsection{L'assioma di completezza}

\begin{assioma}[Completezza]\label{ass:completezza}\index{assioma di completezza}
Siano $A$ e $B$ due insiemi non vuoti di numeri reali con la proprietà che
\[
  a\le b \qquad\text{comunque si scelgano } a\in A \text{ e } b\in B.
\]
Allora esiste almeno un numero reale $c$ tale che
\[
  a\le c\le b \qquad\text{qualunque siano } a\in A \text{ e } b\in B.
\]
\end{assioma}

\begin{marginfigure}
\centering
\begin{tikzpicture}[scale=0.9]
  \draw[-Latex, grigio] (-0.3,0) -- (5.2,0);
  \foreach \x in {0.2,0.6,1.0,1.5,1.9,2.2} \fill[blu] (\x,0) circle (1.6pt);
  \foreach \x in {3.0,3.3,3.8,4.3,4.8} \fill[rame] (\x,0) circle (1.6pt);
  \draw[blu, thick, decorate, decoration={brace, amplitude=4pt, mirror}] (0.1,-0.25) -- (2.3,-0.25)
    node[midway, below=5pt, font=\scriptsize] {$A$};
  \draw[rame, thick, decorate, decoration={brace, amplitude=4pt, mirror}] (2.9,-0.25) -- (4.9,-0.25)
    node[midway, below=5pt, font=\scriptsize] {$B$};
  \draw[verde, thick] (2.6,-0.15) -- (2.6,0.15) node[above, font=\scriptsize] {$c$};
\end{tikzpicture}
\caption{L'assioma di completezza: se ogni elemento di $A$ precede ogni elemento di $B$,
esiste un «separatore» $c$.}
\label{fig:completezza}
\end{marginfigure}

A prima vista l'assioma di completezza può sembrare ovvio, e il lettore potrebbe essere
tentato di «dimostrarlo» così: «se tutti i numeri dell'insieme $A$ sono minori o uguali
di tutti i numeri dell'insieme $B$, allora esisterà certamente un numero $c$
intermedio tra $A$ e $B$; basterà infatti scegliere come $c$ il più grande elemento di
$A$, oppure il più piccolo elemento di $B$».

\begin{attenzione}[La frase precedente è sbagliata!]
Non tutti gli insiemi numerici hanno un più grande o un più piccolo elemento. Ad
esempio, l'insieme
\[
  B=\set{1,\ \tfrac12,\ \tfrac13,\ \tfrac14,\ \dots,\ \tfrac1n,\ \dots},
\]
rappresentato sulla retta nella \cref{fig:insieme-1-su-n}, ha il più grande elemento,
che è $1$, ma \emph{non} ha il più piccolo elemento. Potremmo essere tentati di dire
che lo zero è il più piccolo elemento di $B$; ma lo zero \emph{non è un elemento di
$B$}! Infatti $0\neq 1/n$ qualunque sia $n$, perché una frazione è nulla se e soltanto
se il suo numeratore è nullo. Gli elementi di $B$ si accumulano verso lo zero senza
mai raggiungerlo, e per ogni elemento $1/n$ di $B$ ce n'è uno più piccolo,
$1/(n+1)$.
\end{attenzione}

\begin{figure}[htb]
\centering
\begin{tikzpicture}[scale=1.0]
  \draw[-Latex, grigio] (-0.5,0) -- (11.3,0) node[below left, font=\scriptsize] {};
  \foreach \n in {1,...,60} {\fill[blu] ({10/\n},0) circle (1.3pt);}
  \draw[rosso, thick] (0,0) circle (2.2pt);
  \node[below, font=\scriptsize] at (0,-0.1) {$0$};
  \foreach \n/\lab in {1/1, 2/\tfrac12, 3/\tfrac13, 4/\tfrac14, 5/\tfrac15, 8/\tfrac18} {
    \draw[grigio] ({10/\n},0.08) -- ({10/\n},-0.08);
    \node[below, font=\scriptsize] at ({10/\n},-0.1) {$\lab$};}
\end{tikzpicture}
\caption{L'insieme $B=\{1/n : n\in\N\}$ ha massimo $1$ ma non ha minimo: i suoi
elementi si addensano verso $0$, che però (cerchietto rosso) non appartiene a $B$.}
\label{fig:insieme-1-su-n}
\end{figure}

L'assioma di completezza è, in effetti, un assioma molto più profondo di quanto possa
sembrare a prima vista. Come mostreremo nel \cref{sec:nzq}, soltanto tramite l'assioma di
completezza è possibile distinguere l'insieme dei numeri rappresentabili sotto forma di
frazione, i \emph{numeri razionali}, dall'insieme dei numeri reali: i razionali
soddisfano tutti gli assiomi (A1)--(A6) e (O1)--(O4), ma \emph{non} l'assioma di
completezza. In termini geometrici, l'assioma dice che la retta reale \emph{non ha
buchi}: se la tagliamo in due parti, una a sinistra e una a destra, il punto di taglio
è un numero reale.

\begin{idea}[Un'immagine fisica della completezza]
Si pensi a un'asta graduata e a un cursore che scorre su di essa. Se spingiamo il
cursore verso destra fin dove possiamo, con la regola «puoi avanzare finché resti a
sinistra di tutti i punti di $B$», il cursore si fermerà in una posizione precisa $c$:
non può «cadere in un buco» della retta. L'assioma di completezza garantisce che tale
posizione esista. Sui numeri razionali, invece, il cursore potrebbe trovarsi a dover
scegliere tra posizioni razionali sempre più vicine a un punto che razionale non è,
come $\sqrt2$, e non avrebbe dove fermarsi.
\end{idea}

\begin{python}[Verifica numerica delle proprietà di campo]
# Gli assiomi (A1)-(A6) si possono "toccare con mano" con numeri
# esatti: Fraction rappresenta i razionali senza arrotondamenti.
from fractions import Fraction as F

a, b, c = F(3, 4), F(-5, 6), F(7, 2)
print((a + b) + c == a + (b + c))      # (A1) associativa (somma)
print((a * b) * c == a * (b * c))      # (A1) associativa (prodotto)
print(a + b == b + a, a * b == b * a)  # (A2) commutativa
print(a * (b + c) == a * b + a * c)    # (A3) distributiva
print(a + 0 == a, a * 1 == a)          # (A4) elementi neutri
print(a + (-a) == 0)                   # (A5) opposto
print(a * (1 / a) == 1)                # (A6) inverso (a != 0)
\end{python}
\begin{risultato}
True
True
True True
True
True True
True
True
\end{risultato}

\begin{attenzione}[I numeri in virgola mobile non sono numeri reali]
Se nel codice precedente si sostituisce \texttt{Fraction} con i normali numeri
decimali (\texttt{float}), alcune uguaglianze possono fallire: ad esempio in Python
\texttt{0.1 + 0.2 == 0.3} restituisce \texttt{False}, e la proprietà associativa
\texttt{(0.1 + 0.2) + 0.3 == 0.1 + (0.2 + 0.3)} è anch'essa \texttt{False}. I
\texttt{float} sono un insieme \emph{finito} di numeri che approssima $\R$: ottimi per
calcolare, ma non obbediscono agli assiomi. Questa distinzione tra il modello matematico
(i numeri reali) e la sua rappresentazione al calcolatore è una delle prime cose che un
ingegnere deve avere chiara.
\end{attenzione}
